\documentclass[12pt]{article}
%\usepackage{fullpage} IF 
\usepackage[a4paper, margin=1in]{geometry} % this one is less ugly 

\usepackage{amsmath}
\usepackage{mathtools}
\usepackage{setspace}
\usepackage{pgfgantt}

\doublespacing

\begin{document}

\title{Project Initiation Document: Remaining time prediction}
\author{Group 1: Lev Fedotov, Mariia Vashchenko, Kiarash Kowsari}
\date{October 2026}
\maketitle
\pagebreak
\tableofcontents
%--INTRODUCTION--%
\section{Introduction}
In this project, we aim to implement software for remaining time prediction for the [A1 GmbH], which is a consulting
company wanting to optimize their service. This initiation document contains
every detail for initiating this project. \\
In chapter \textbf{2}, the business cases, for which the product will be developed, will
be presented. The nancial importance of our product for the client will be
explained and two scenarios where our product can be used to improve the
service provided by the A1 GmbH will be listed. \\ 
In chapter \textbf{3}, to discuss the feasibility of our product, the concepts and tools
that will be used to implement our product will be introduced. Two use cases
will also be given to show the applicability of our product. \\
In chapter \textbf{4}, a GANTT chart outlining the time frame of our project, and its
explanation will be provided. \\ 
In chapter \textbf{5}, tools that will be used in this project and the reason why we
choose to use these tools will be given.
In the \textbf{6}th chapter, an overview of our development team will be provided, which
contains information about our background, academic interests, strengths, and
weaknesses as well as the roles assigned to each member of our team.
In the last chapter, we will  
\section{Business Case}
\subsection{Client}
\section{Feasibility Study}
\subsection{Theoretical and Technical Feasibility}
\subsection{Use Cases}
\subsubsection{Python}
\section{Project charter}

\definecolor{barblue}{RGB}{153,204,254}
\definecolor{groupblue}{RGB}{51,102,254}
\definecolor{linkred}{RGB}{165,0,33}
\setganttlinklabel{s-s}{START-TO-START}
\setganttlinklabel{f-s}{FINISH-TO-START}
\setganttlinklabel{f-f}{FINISH-TO-FINISH}
\begin{ganttchart}[
    canvas/.append style={fill=none, draw=black!5, line width=.75pt},
    hgrid style/.style={draw=black!5, line width=.75pt},
    vgrid={*1{draw=black!5, line width=.75pt}},
    today=7,
    today rule/.style={
      draw=black!64,
      dash pattern=on 3.5pt off 4.5pt,
      line width=1.5pt
    },
    today label font=\small\bfseries,
    title/.style={draw=none, fill=none},
    title label font=\bfseries\footnotesize,
    title label node/.append style={below=7pt},
    include title in canvas=false,
    bar label font=\mdseries\small\color{black!70},
    bar label node/.append style={left=2cm},
    bar/.append style={draw=none, fill=black!63},
    bar incomplete/.append style={fill=barblue},
    bar progress label font=\mdseries\footnotesize\color{black!70},
    group incomplete/.append style={fill=groupblue},
    group left shift=0,
    group right shift=0,
    group height=.5,
    group peaks tip position=0,
    group label node/.append style={left=.6cm},
    group progress label font=\bfseries\small,
    link/.style={-latex, line width=1.5pt, linkred},
    link label font=\scriptsize\bfseries,
    link label node/.append style={below left=-2pt and 0pt}
  ]{1}{13}
  \gantttitle[
    title label node/.append style={below left=7pt and -3pt}
  ]{WEEKS:\quad1}{1}
  \gantttitlelist{2,...,13}{1} \\
  \ganttgroup[progress=57]{WBS 1 Summary Element 1}{1}{10} \\
  \ganttbar[
    progress=75,
    name=WBS1A
  ]{\textbf{WBS 1.1} Activity A}{1}{8} \\
  \ganttbar[
    progress=67,
    name=WBS1B
  ]{\textbf{WBS 1.2} Activity B}{1}{3} \\
  \ganttbar[
    progress=50,
    name=WBS1C
  ]{\textbf{WBS 1.3} Activity C}{4}{10} \\
  \ganttbar[
    progress=0,
    name=WBS1D
  ]{\textbf{WBS 1.4} Activity D}{4}{10} \\[grid]
  \ganttgroup[progress=0]{WBS 2 Summary Element 2}{4}{10} \\
  \ganttbar[progress=0]{\textbf{WBS 2.1} Activity E}{4}{5} \\
  \ganttbar[progress=0]{\textbf{WBS 2.2} Activity F}{6}{8} \\
  \ganttbar[progress=0]{\textbf{WBS 2.3} Activity G}{9}{10}
  \ganttlink[link type=s-s]{WBS1A}{WBS1B}
  \ganttlink[link type=f-s]{WBS1B}{WBS1C}
  \ganttlink[
    link type=f-f,
    link label node/.append style=left
  ]{WBS1C}{WBS1D}
\end{ganttchart}

\begin{itemize}
    \item MILESTONE 2
    \item MILESTONE 3
    \item MILESTONE 4
    \item MILESTONE 5
\end{itemize}


\section{Tools}
\subsection{Gitlab of the RWTH Aachen Univeristy}
RWTH Aachen is a member university of NFDI4Ing and the students have
the admission to use DFN AAI Single Sign-On to create projects and groups.
We use GitLab of the RWTH Aachen University to manage our own Gitrepository so that each team member has access to our project and to
increase the quality of the code. Besides, this website has other functions as
an issue-tracking-system, a wiki, and an option to review codes. The
advantage of utilizing GitLab is that external users can authenticate via
GitHub, however, they do not have the access to create groups and proj
\section{Team}
\subsection{Personal description}
\subsubsection{Lev Fedotov}
I am a B.Sc Computer Science student at RWTH Aachen University in my 5th semester, with my electives mostly being in math and psychology. Right now I am also employed as a research assistant at Fraunhofer Research Institution for Battery Cell Production FFB, primarily doing full stack development for a manufacturing execution system. Although most of my practical experience is in web development, my academic interests remain quite broad so I am eager to explore subfields I have not yet engaged with. I haven't engaged before  . In this practicum, my aim is to contribute to the project’s success while learning to apply my existing theoretical knowledge in a more practical, real-world like setting, which could help me in other areas of CS and more specifically in my current niche of digital production planning.

\subsubsection{Mariia Vaschenko}
I am a 5th-semester Bachelor's student in Computer Science at RWTH Aachen University, with a particular interest in Machine Learning and Neural Networks. My primary programming language is Python, but I also have experience in Java and C++. In a recently finished Computer Vision project, I used NumPy for image preprocessing and PyTorch to implement, train, and evaluate a neural network. I am also working on a market mispricing prediction project involving REST APIs and data preprocessing.
Currently, I work as a teaching assistant at the Institute for AI Methodology, where I conduct weekly tutorials and assess students' assignments. This experience has strengthened both my technical knowledge and my ability to explain complex topics and communicate effectively.
Although I have experience in machine learning, my knowledge of process mining is still limited. I hope to contribute my experience in data preprocessing and predictive modelling to the development and evaluation of regression and tree-based models while gaining a deeper understanding of event-log analysis and collaborative software development.

\subsubsection{Kiarash Kowsari}
I am a B.Sc. Computer Science student in my 5th semester at RWTH Aachen University. My electives focus on process mining and computer vision, and alongside them I am following courses on machine learning and agentic AI. For about a year I have worked as a backend developer at Kisters AG, mainly on infrastructure, using Spring Boot, Go, and DevOps tooling such as Docker, Kubernetes, Jenkins, Grafana, and OpenTelemetry. There I work in an agile Scrum environment. I have also worked with the process mining tools ProM, Celonis, and RapidMiner, and I am comfortable with data analysis and basic machine learning in Python. In addition, I serve as a team leader in a data science club, where I have gained experience in leading teams. My main strengths are learning new skills quickly and organizing a team's work. My interests lie in AI for infrastructure, agentic AI, and explainable AI, and I like to apply new approaches in industry. On the other hand, I have not yet built self-created projects or written academic papers to showcase my skills, and I still have much to learn in my fields. This lab is therefore an excellent opportunity for me: it combines my interest in process mining with applied machine learning, and remaining time prediction is a task I could use in industry in the near future. I am looking forward to working with my team to deliver a complete project.

\subsection{Roles}
[DISCUSS IN THE MEETING ON 2.11.26]
%--BIBLIOGRAPHY--%
\section{Phase review}
\subsection{Lev Fedotov}
It was very stupid of me to miss the first team specific meeting // TODO REWRITE
[please write after   the meeting on 2.11.26]

\subsection{Mariia Vaschenko}
[please write after   the meeting on 2.11.26]

\subsection{Kiarash Kowsari}
[please write after  the meeting on 2.11.26]

\begin{thebibliography}{99}

\bibitem{SampleBibLabel} Team0, "Title",  Journal, Volume, Pages, Year. 

\end{thebibliography}

\end{document}
